\chapter{2013年天津卷（文）}

\section{单选题}

\begin{problem}
已知集合 \(A = \{x \in \R \mid |x| \leqslant 2\}\)，\(B = \{x \in \R \mid x \leqslant 1\}\)，则 \(A \cap B =\)（\quad）

\choices
    {\((-\infty, 2]\)}
    {\([1,2]\)}
    {\([-2, 2]\)}
    {\([-2, 1]\)}
\end{problem}

\begin{answer}
D
\end{answer}

\begin{solution}
由 \(|x|\leqslant 2\) 得 \(-2\leqslant x\leqslant 2\)，再与 \(x\leqslant 1\) 取公共部分，得
\(A\cap B=[-2,1]\)，故选 D.
\end{solution}

\begin{problem}
设变量 \(x\)，\(y\) 满足约束条件 \(\begin{cases}
3x + y - 6 \geqslant 0,\\
x - y - 2 \leqslant 0,\\
y - 3 \leqslant 0,
\end{cases}\) 则目标函数 \(z = y - 2x\) 的最小值为（\quad）

\choices
    {\(-7\)}
    {\(-4\)}
    {1}
    {2}
\end{problem}

\begin{answer}
A
\end{answer}

\begin{solution}
约束条件等价于 \(y\geqslant 6-3x\)，\(y\geqslant x-2\)，\(y\leqslant 3\). 三条边界直线的交点分别为：
\(y=3\) 与 \(y=6-3x\) 交于 \((1,3)\)，\(y=3\) 与 \(y=x-2\) 交于 \((5,3)\)，而 \(y=6-3x\) 与 \(y=x-2\) 交于 \((2,0)\). 因此可行域的三个顶点为
\((1,3)\)，\((5,3)\)，\((2,0)\). 目标函数在这三个顶点处分别为
\(1\)，\(-7\)，\(-4\). 因此 \(z\) 的最小值为 \(-7\)，故选 A.
\end{solution}

\begin{problem}
阅读如图所示的程序框图，运行相应的程序，则输出 \(n\) 的值为（\quad）

\bitmapfigure[max width=0.42\linewidth]{img/2013/tianjin_liberal/q3_fig1.png}

\choices
    {7}
    {6}
    {5}
    {4}
\end{problem}

\begin{answer}
D
\end{answer}

\begin{solution}
初始时 \(n=1\)，\(S=0\). 每次先作 \(S\leftarrow S+(-1)^n n\)，再判断是否有 \(S\geqslant 2\). 逐次计算如下：
当 \(n=1\) 时，\(S=-1\)；当 \(n=2\) 时，\(S=1\)；当 \(n=3\) 时，\(S=-2\)；当 \(n=4\) 时，\(S=2\). 此时满足判断条件，输出 \(n=4\)，故选 D.
\end{solution}

\begin{problem}
设 \(a\)，\(b \in \R\)，则“\((a - b) \cdot a^2 < 0\)”是“\(a < b\)”的（\quad）

\choices
    {充分而不必要条件}
    {必要而不充分条件}
    {充要条件}
    {既不充分也不必要条件}
\end{problem}

\begin{answer}
A
\end{answer}

\begin{solution}
若 \((a-b)a^2<0\)，则必有 \(a\neq 0\)，从而 \(a^2>0\). 因此只能是 \(a-b<0\)，即 \(a<b\)，所以该条件是 \(a<b\) 的充分条件.

反之，取 \(a=0\)，\(b=1\)，虽有 \(a<b\)，但 \((a-b)a^2=0\) 不小于 \(0\)，故不是必要条件. 因此选 A.
\end{solution}

\begin{problem}
已知过点 \(P(2,2)\) 的直线与 \((x - 1)^2 + y^2 = 5\) 相切，且与直线 \(ax - y + 1 = 0\) 垂直，则 \(a =\)（\quad）

\choices
    {\(-\frac{1}{2}\)}
    {1}
    {2}
    {\(\frac{1}{2}\)}
\end{problem}

\begin{answer}
C
\end{answer}

\begin{solution}
圆心为 \(C(1,0)\)，且 \(CP=\sqrt{(2-1)^2+2^2}=\sqrt{5}\)，所以 \(P\) 在圆上. 过圆上点 \(P\) 的切线与半径 \(CP\) 垂直.

半径 \(CP\) 的斜率为 \(2\)，故切线斜率为 \(-\frac{1}{2}\). 直线 \(ax-y+1=0\) 的斜率为 \(a\)，它与切线垂直，所以
\(a\cdot\left(-\frac{1}{2}\right)=-1\)，解得 \(a=2\)，故选 C.
\end{solution}

\begin{problem}
函数 \( f(x) = \sin \left(2x - \frac{\pi}{4}\right) \) 在区间 \(\left[0, \frac{\pi}{2}\right]\) 上的最小值为（\quad）

\choices
    {\(-1\)}
    {\(-\frac{\sqrt{2}}{2}\)}
    {\(\frac{\sqrt{2}}{2}\)}
    {0}
\end{problem}

\begin{answer}
B
\end{answer}

\begin{solution}
令 \(t=2x-\frac{\pi}{4}\). 当 \(x\in\left[0,\frac{\pi}{2}\right]\) 时，
\(t\in\left[-\frac{\pi}{4},\frac{3\pi}{4}\right]\). 在此区间内，
\(\sin t\) 的导数为 \(\cos t\)，唯一驻点为 \(t=\frac{\pi}{2}\)，且 \(\sin t\) 在
\(\left[-\frac{\pi}{4},\frac{\pi}{2}\right]\) 上递增、在
\(\left[\frac{\pi}{2},\frac{3\pi}{4}\right]\) 上递减. 因此比较两个端点和驻点的函数值，得
\[
\min f(x)=\sin\left(-\frac{\pi}{4}\right)=-\frac{\sqrt{2}}{2}
\]
故选 B.
\end{solution}

\begin{problem}
已知函数 \( f(x) \) 是定义在 \( \R \) 上的偶函数，且在区间 \( [0, +\infty) \) 上单调递增. 若实数 \( a \) 满足 \( f(\log_{2} a) + f(\log_{\frac{1}{2}} a) \leqslant 2f(1) \)，则 \( a \) 的取值范围是（\quad）

\choices
    {\([1,2]\)}
    {\(\left(0, \frac{1}{2}\right]\)}
    {\(\left[\frac{1}{2}, 2\right]\)}
    {\((0,2]\)}
\end{problem}

\begin{answer}
C
\end{answer}

\begin{solution}
由于 \(a>0\)，令 \(t=\log_{2}a\)，则 \(\log_{\frac{1}{2}}a=-t\). 利用 \(f\) 为偶函数，原不等式化为
\(2f(t)\leqslant 2f(1)\)，即 \(f(|t|)\leqslant f(1)\).

由 \(f\) 在 \([0,+\infty)\) 上单调递增，得 \(|t|\leqslant 1\)，于是
\(-1\leqslant\log_{2}a\leqslant 1\). 因而 \(\frac{1}{2}\leqslant a\leqslant 2\)，故选 C.
\end{solution}

\begin{problem}
设函数 \( f(x) = \e^{x} + x - 2 \)，\( g(x) = \ln x + x^2 - 3 \). 若实数 \(a\)，\(b\) 满足 \( f(a) = 0 \)，\( g(b) = 0 \)，则（\quad）

\choices
    {\(g(a) < 0 < f(b)\)}
    {\(f(b) < 0 < g(a)\)}
    {\(0 < g(a) < f(b)\)}
    {\(f(b) < g(a) < 0\)}
\end{problem}

\begin{answer}
A
\end{answer}

\begin{solution}
因为 \(f'(x)=\e^x+1>0\)，所以 \(f\) 在 \(\R\) 上单调递增. 又 \(f(0)=-1<0\)，\(f(1)=\e-1>0\)，故由 \(f(a)=0\) 得 \(0<a<1\).

因为 \(g'(x)=\frac{1}{x}+2x>0\)（\(x>0\)），所以 \(g\) 在 \((0,+\infty)\) 上单调递增. 又 \(g(1)=-2<0\)，\(g(2)=\ln 2+1>0\)，故由 \(g(b)=0\) 得 \(1<b<2\).

由 \(0<a<1\) 得 \(\ln a<0\) 且 \(a^2<1\)，所以 \(g(a)=\ln a+a^2-3<0\). 由 \(b>1\) 及 \(f\) 的单调性，得 \(f(b)>f(1)>0\). 因此 \(g(a)<0<f(b)\)，故选 A.
\end{solution}

\section{填空题}

\begin{problem}
\(\i\) 是虚数单位，复数 \( (3+\i)(1-2\i)= \)\fillinblank{}.
\end{problem}

\begin{answer}
\(5-5\i\)
\end{answer}

\begin{solution}
利用 \(\i^2=-1\)，有
\[
(3+\i)(1-2\i)=3-6\i+\i-2\i^2=3-5\i+2=5-5\i
\]
所以填 \(5-5\i\).
\end{solution}

\begin{problem}
已知一个正方体的所有顶点在一个球面上，若球的体积为 \(\frac{9\pi}{2}\)，则正方体的棱长为\fillinblank{}.
\end{problem}

\begin{answer}
\(\sqrt{3}\)
\end{answer}

\begin{solution}
设球的半径为 \(R\). 由球的体积得
\(\frac{4}{3}\pi R^3=\frac{9\pi}{2}\)，所以 \(R=\frac{3}{2}\).

正方体的体对角线等于外接球的直径 \(2R=3\). 设正方体棱长为 \(l\)，则
\(l\sqrt{3}=3\)，解得 \(l=\sqrt{3}\). 所以填 \(\sqrt{3}\).
\end{solution}

\begin{problem}
已知抛物线 \( y^{2} = 8x \) 的准线过双曲线 \( \frac{x^2}{a^2} - \frac{y^2}{b^2} = 1  \) （\(a > 0\)，\(b > 0\)）的一个焦点，且双曲线的离心率为 2，则该双曲线的方程为\fillinblank{}.
\end{problem}

\begin{answer}
\(x^2-\frac{y^2}{3}=1\)
\end{answer}

\begin{solution}
抛物线 \(y^2=8x\) 可写为 \(y^2=4px\)，故 \(p=2\)，其准线为 \(x=-2\). 设双曲线焦半距为 \(c\)，则一个焦点在准线上给出 \(c=2\).

由离心率 \(e=\frac{c}{a}=2\)，得 \(a=1\). 又 \(c^2=a^2+b^2\)，所以 \(b^2=4-1=3\). 因此双曲线方程为
\(\frac{x^2}{1^2}-\frac{y^2}{3}=1\)，即 \(x^2-\frac{y^2}{3}=1\).
\end{solution}

\begin{problem}
在平行四边形 \(ABCD\) 中，\(AD = 1\)，\(\angle BAD = 60^{\circ}\)，\(E\) 为 \(CD\) 的中点. 若 \(\overrightarrow{AC} \cdot \overrightarrow{BE} = 1\)，则 \(AB\) 的长为\fillinblank{}.
\end{problem}

\begin{answer}
\(\frac{1}{2}\)
\end{answer}

\begin{solution}
设 \(\overrightarrow{AB}=\bs{u}\)，\(\overrightarrow{AD}=\bs{v}\)，则
\(|\bs{v}|=1\)，\(|\bs{u}|=AB\)，且 \(\bs{u}\cdot\bs{v}=\frac{1}{2}|\bs{u}|\).

取 \(A\) 为原点，则 \(C\) 的位置向量为 \(\bs{u}+\bs{v}\)，而 \(E\) 为 \(CD\) 的中点，所以
\(\overrightarrow{BE}=\bs{v}-\frac{1}{2}\bs{u}\). 因而
\[
\overrightarrow{AC}\cdot\overrightarrow{BE}
=(\bs{u}+\bs{v})\cdot\left(\bs{v}-\frac{1}{2}\bs{u}\right)
=1+\frac{1}{4}|\bs{u}|-\frac{1}{2}|\bs{u}|^2
\]
令 \(t=|\bs{u}|>0\)，由题设得
\(1+\frac{t}{4}-\frac{t^2}{2}=1\)，即 \(t(1-2t)=0\). 因为 \(t>0\)，所以 \(AB=t=\frac{1}{2}\).
\end{solution}

\begin{problem}
如图，在圆内接梯形 \(ABCD\) 中，\(AB \parallel DC\)，过点 \(A\) 作圆的切线与 \(CB\) 的延长线交于点 \(E\). 若 \(AB = AD = 5\)，\(BE = 4\)，则弦 \(BD\) 的长为\fillinblank{}.

\bitmapfigure[max width=0.66\linewidth]{img/2013/tianjin_liberal/q13_fig1.png}
\end{problem}

\begin{answer}
\(\frac{15}{2}\)
\end{answer}

\begin{solution}
因为圆内接梯形的两底平行，所以它是等腰梯形，从而 \(BC=AD=5\). 又 \(C\)，\(B\)，\(E\) 共线且 \(BE=4\)，故 \(CE=9\).

由切线与割线定理，\(EA^2=EB\cdot EC=4\cdot 9=36\)，所以 \(EA=6\). 在 \(\triangle ABE\) 中，由余弦定理
\[
\cos\angle ABE=\frac{AB^2+BE^2-AE^2}{2\cdot AB\cdot BE}=\frac{25+16-36}{40}=\frac{1}{8}
\]
由于 \(BE\) 是 \(BC\) 的反向延长线，\(\angle ABE=\pi-\angle ABC\)，故 \(\cos\angle ABC=-\frac{1}{8}\).

又由 \(AB\parallel DC\) 及四边形 \(ABCD\) 圆内接，有 \(\angle DAB=\angle ABC\)，所以 \(\cos\angle DAB=-\frac{1}{8}\). 在等腰三角形 \(\triangle ABD\) 中，\(AB=AD=5\)，从而
\[
BD^2=AB^2+AD^2-2AB\cdot AD\cos\angle DAB
=25+25-50\left(-\frac{1}{8}\right)=\frac{225}{4}
\]
因此 \(BD=\frac{15}{2}\).
\end{solution}

\begin{problem}
设 \(a + b = 2\)，\(b > 0\)，则 \(\frac{1}{2|a|} + \frac{|a|}{b}\) 的最小值为\fillinblank{}.
\end{problem}

\begin{answer}
\(\frac{3}{4}\)
\end{answer}

\begin{solution}
因为分母中有 \(|a|\)，所以 \(a\neq 0\). 分两种情况讨论.

当 \(a<0\) 时，令 \(x=|a|=-a>0\)，则 \(b=x+2\)，且
\[
\frac{1}{2|a|}+\frac{|a|}{b}-\frac{3}{4}
=\frac{1}{2x}+\frac{x}{x+2}-\frac{3}{4}
=\frac{(x-2)^2}{4x(x+2)}\geqslant 0
\]
等号在 \(x=2\)，即 \(a=-2\)，\(b=4\) 时取得.

当 \(a>0\) 时，令 \(x=a\)，则 \(0<x<2\)，\(b=2-x\)，并且
\[
\frac{1}{2|a|}+\frac{|a|}{b}-\frac{5}{4}
=\frac{1}{2x}+\frac{x}{2-x}-\frac{5}{4}
=\frac{(3x-2)^2}{4x(2-x)}\geqslant 0
\]
此时原式最小为 \(\frac{5}{4}>\frac{3}{4}\). 综上，原式的最小值为 \(\frac{3}{4}\).
\end{solution}

\section{解答题}

\begin{problem}
某产品的三个质量指标分别为 \(x\)，\(y\)，\(z\)，用综合指标 \(S = x + y + z\) 评价该产品的等级. 若 \(S \leqslant 4\)，则该产品为一等品. 现从一批该产品中，随机抽取 \(10\) 件产品作为样本，其质量指标列表如下：

\begin{center}
\begingroup
\small
\renewcommand{\arraystretch}{1.2}
\begin{tabular}{|c|c|c|c|c|c|}
\hline
产品编号 & \(A_1\) & \(A_2\) & \(A_3\) & \(A_4\) & \(A_5\) \\
\hline
质量指标 \((x,y,z)\) & \((1,1,2)\) & \((2,1,1)\) & \((2,2,2)\) & \((1,1,1)\) & \((1,2,1)\) \\
\hline
产品编号 & \(A_6\) & \(A_7\) & \(A_8\) & \(A_9\) & \(A_{10}\) \\
\hline
质量指标 \((x,y,z)\) & \((1,2,2)\) & \((2,1,1)\) & \((2,2,1)\) & \((1,1,1)\) & \((2,1,2)\) \\
\hline
\end{tabular}
\endgroup
\end{center}

\begin{enumerate}
\item 利用上表提供的样本数据估计出该批产品的一等品率；
\item 在该样本的一等品中，随机抽取 \(2\) 件产品.
\begin{circlelist}
\item 用产品编号列出所有可能的结果；
\item 设事件 \(B\) 为“在取出的 \(2\) 件产品中，每件产品的综合指标 \(S\) 都等于 \(4\)”，求事件 \(B\) 发生的概率.
\end{circlelist}
\end{enumerate}
\end{problem}

\begin{solution}
\begin{enumerate}
\item 逐件计算综合指标 \(S=x+y+z\)，得到
\(S_{A_1}=4\)，\(S_{A_2}=4\)，\(S_{A_3}=6\)，\(S_{A_4}=3\)，\(S_{A_5}=4\)，\(S_{A_6}=5\)，\(S_{A_7}=4\)，\(S_{A_8}=5\)，\(S_{A_9}=3\)，\(S_{A_{10}}=5\).
其中 \(S\leqslant 4\) 的产品为 \(A_1\)，\(A_2\)，\(A_4\)，\(A_5\)，\(A_7\)，\(A_9\)，共有 \(6\) 件. 因此用样本一等品数估计该批产品的一等品率为
\(\frac{6}{10}=0.6\).

\item 在样本的一等品中共有 \(6\) 件，任取 \(2\) 件且不计顺序，所有可能结果为
\(\{A_1,A_2\}\)，\(\{A_1,A_4\}\)，\(\{A_1,A_5\}\)，\(\{A_1,A_7\}\)，\(\{A_1,A_9\}\)，\(\{A_2,A_4\}\)，\(\{A_2,A_5\}\)，\(\{A_2,A_7\}\)，\(\{A_2,A_9\}\)，\(\{A_4,A_5\}\)，\(\{A_4,A_7\}\)，\(\{A_4,A_9\}\)，\(\{A_5,A_7\}\)，\(\{A_5,A_9\}\)，\(\{A_7,A_9\}\).

 \begin{enumerate}
 \item 以上 \(15\) 个无序二元集合就是所有可能结果.

 \item 综合指标等于 \(4\) 的一等品为 \(A_1\)，\(A_2\)，\(A_5\)，\(A_7\)，所以事件 \(B\) 包含 \(\mathrm C_4^2=6\) 个结果，而所有结果共有 \(\mathrm C_6^2=15\) 个. 因此
 \(P(B)=\frac{\mathrm C_4^2}{\mathrm C_6^2}=\frac{6}{15}=\frac{2}{5}\).
 \end{enumerate}
\end{enumerate}
\end{solution}

\begin{problem}
在 \(\triangle ABC\) 中，内角 \(A\)，\(B\)，\(C\) 所对的边分别是 \(a\)，\(b\)，\(c\)，已知 \(b \sin A = 3c \sin B\)，\(a = 3\)，\(\cos B = \frac{2}{3}\).
\begin{enumerate}
\item 求 \(b\) 的值；
\item 求 \( \sin\left(2B-\frac{\pi}{3}\right) \) 的值.
\end{enumerate}
\end{problem}

\begin{solution}
\begin{enumerate}
\item 由正弦定理，\(\frac{a}{\sin A}=\frac{b}{\sin B}=\frac{c}{\sin C}\)，所以 \(b\sin A=a\sin B\). 结合题设 \(b\sin A=3c\sin B\)，且 \(\sin B>0\)，得 \(a=3c\). 因为 \(a=3\)，所以 \(c=1\).

由余弦定理及 \(\cos B=\frac{2}{3}\)，有
\[
b^2=a^2+c^2-2ac\cos B=3^2+1^2-2\cdot3\cdot1\cdot\frac{2}{3}=6
\]
由于 \(b>0\)，故 \(b=\sqrt{6}\).

\item 因为 \(B\) 是三角形内角，所以 \(\sin B=\sqrt{1-\cos^2B}=\frac{\sqrt{5}}{3}\). 于是
\(\sin 2B=2\sin B\cos B=\frac{4\sqrt{5}}{9}\)，\(\cos 2B=2\cos^2B-1=-\frac{1}{9}\). 因此
\[
\sin\left(2B-\frac{\pi}{3}\right)
=\sin 2B\cos\frac{\pi}{3}-\cos 2B\sin\frac{\pi}{3}
=\frac{4\sqrt{5}}{9}\cdot\frac{1}{2}+\frac{1}{9}\cdot\frac{\sqrt{3}}{2}
=\frac{4\sqrt{5}+\sqrt{3}}{18}
\]
\end{enumerate}
\end{solution}

\begin{problem}
如图，三棱柱 \(ABC - A_{1}B_{1}C_{1}\) 中，侧棱 \(A_{1}A \perp\) 底面 \(ABC\)，且各棱长均相等，\(D\)，\(E\)，\(F\) 分别为棱 \(AB\)，\(BC\)，\(A_{1}C_{1}\) 的中点.
\begin{enumerate}
\item 证明 \(EF \parallel\) 平面 \(A_{1}CD\)；
\item 证明平面 \(A_{1}CD \perp\) 平面 \(A_{1}ABB_{1}\)；
\item 求直线 \(BC\) 与平面 \(A_{1}CD\) 所成角的正弦值.
\end{enumerate}

\bitmapfigure[max width=0.46\linewidth]{img/2013/tianjin_liberal/q17_fig1.png}
\end{problem}

\begin{solution}
因三棱柱的各棱长均相等，底面 \(ABC\) 为边长相等的等边三角形. 设各棱长为 \(l>0\)，建立空间直角坐标系，使
\(A(0,0,0)\)，\(B(l,0,0)\)，\(C\left(\frac l2,\frac{\sqrt{3}l}{2},0\right)\)，\(A_1(0,0,l)\).
则
\(D\left(\frac l2,0,0\right)\)，\(E\left(\frac{3l}{4},\frac{\sqrt{3}l}{4},0\right)\)，\(F\left(\frac l4,\frac{\sqrt{3}l}{4},l\right)\).

\begin{enumerate}
\item 平面 \(A_1CD\) 经过 \(A_1\)，\(C\)，\(D\)，其方程为 \(2x+z=l\)，故可取法向量 \(\bs{n}=(2,0,1)\). 又
\(\overrightarrow{EF}=\left(-\frac l2,0,l\right)\)，从而 \(\overrightarrow{EF}\cdot\bs{n}=0\). 另外点 \(E\) 不在平面 \(A_1CD\) 上，因此 \(EF\parallel\) 平面 \(A_1CD\).

\item 平面 \(A_1ABB_1\) 的方程为 \(y=0\)，可取法向量 \(\bs{m}=(0,1,0)\). 因为 \(\bs{n}\cdot\bs{m}=0\)，所以两平面的法向量垂直，故平面 \(A_1CD\perp\) 平面 \(A_1ABB_1\).

\item 设直线 \(BC\) 与平面 \(A_1CD\) 所成的角为 \(\theta\). 取直线 \(BC\) 的单位方向向量
\(\bs{v}=\left(-\frac12,\frac{\sqrt{3}}{2},0\right)\)，则
\[
\sin\theta=\frac{|\bs{v}\cdot\bs{n}|}{|\bs{v}|\,|\bs{n}|}
=\frac{1}{\sqrt{5}}=\frac{\sqrt{5}}{5}
\]
\end{enumerate}
\end{solution}

\begin{problem}
设椭圆 \(\frac{x^2}{a^2} + \frac{y^2}{b^2} = 1 \) （\(a > b > 0\)）的左焦点为 \(F\)，离心率为 \(\frac{\sqrt{3}}{3}\)，过点 \(F\) 且与 \(x\) 轴垂直的直线被椭圆截得的线段长为 \(\frac{4\sqrt{3}}{3}\).
\begin{enumerate}
\item 求椭圆的方程；
\item 设 \(A\)，\(B\) 分别为椭圆的左，右顶点，过点 \(F\) 且斜率为 \(k\) 的直线与椭圆交于 \(C\)，\(D\) 两点. 若 \(\overrightarrow{AC} \cdot \overrightarrow{DB} + \overrightarrow{AD} \cdot \overrightarrow{CB} = 8\)，求 \(k\) 的值.
\end{enumerate}
\end{problem}

\begin{solution}
\begin{enumerate}
\item 设椭圆焦半距为 \(c\)，则 \(\frac{c}{a}=\frac{\sqrt{3}}{3}\)，即 \(c^2=\frac{a^2}{3}\). 过左焦点的竖直线为 \(x=-c\)，代入椭圆方程，所得弦长为
\[
2b\sqrt{1-\frac{c^2}{a^2}}=2b\sqrt{\frac{2}{3}}=\frac{4\sqrt{3}}{3}
\]
故 \(b^2=2\). 又 \(b^2=a^2-c^2=\frac{2a^2}{3}\)，所以 \(a^2=3\)，进而 \(c=1\). 椭圆方程为
\(\frac{x^2}{3}+\frac{y^2}{2}=1\).

\item 此时 \(A(-\sqrt{3},0)\)，\(B(\sqrt{3},0)\)，\(F(-1,0)\). 设
\(C(x_1,y_1)\)，\(D(x_2,y_2)\)，则直线 \(CD\) 的方程为 \(y=k(x+1)\). 与椭圆联立，得
\[
(2+3k^2)x^2+6k^2x+3k^2-6=0
\]
所以
\(x_1+x_2=-\frac{6k^2}{2+3k^2}\)，\(x_1x_2=\frac{3k^2-6}{2+3k^2}\).
又 \(y_i=k(x_i+1)\)，故
\[
y_1y_2=k^2(x_1+1)(x_2+1)=-\frac{4k^2}{2+3k^2}
\]
由向量坐标表示，
\(\overrightarrow{AC}=\left(x_1+\sqrt{3},y_1\right)\)，\(\overrightarrow{DB}=\left(\sqrt{3}-x_2,-y_2\right)\)
\(\overrightarrow{AD}=\left(x_2+\sqrt{3},y_2\right)\)，\(\overrightarrow{CB}=\left(\sqrt{3}-x_1,-y_1\right)\).
因此题设左端为
\[
\overrightarrow{AC}\cdot\overrightarrow{DB}
+\overrightarrow{AD}\cdot\overrightarrow{CB}
=2\left(3-x_1x_2-y_1y_2\right)
\]
代入 \(x_1x_2\) 和 \(y_1y_2\) 的值，得
\[
2\left(3-\frac{3k^2-6}{2+3k^2}+\frac{4k^2}{2+3k^2}\right)
=\frac{24+20k^2}{2+3k^2}
\]
令其等于 \(8\)，得 \(24+20k^2=16+24k^2\)，即 \(k^2=2\). 因此 \(k=\pm\sqrt{2}\).
\end{enumerate}
\end{solution}

\begin{problem}
已知首项为 \(\frac{3}{2}\) 的等比数列 \(\{a_{n}\}\) 的前 \(n\) 项和为 \(S_{n}\) （\(n \in \mathbf{N}^{*}\)），且 \(-2S_{2}\)，\(S_{3}\)，\(4S_{4}\) 成等差数列.
\begin{enumerate}
\item 求数列 \( \{a_{n}\} \) 的通项公式；
\item 证明 \(S_{n} + \frac{1}{S_{n}} \leqslant \frac{13}{6} \) （\(n \in \mathbf{N}^{*}\)）.
\end{enumerate}
\end{problem}

\begin{solution}
设等比数列的公比为 \(q\)，由等差数列条件得
\[
S_3-(-2S_2)=4S_4-S_3
\]
即 \(S_3+S_2=2S_4\). 代入
\(S_2=\frac{3}{2}(1+q)\)，\(S_3=\frac{3}{2}(1+q+q^2)\)，\(S_4=\frac{3}{2}(1+q+q^2+q^3)\)，得
\(q^2(1+2q)=0\). 等比数列的公比不为零，故 \(q=-\frac{1}{2}\)，从而
\[
a_n=\frac{3}{2}\left(-\frac{1}{2}\right)^{n-1}
\]

又
\[
S_n=\frac{a_1(1-q^n)}{1-q}=1-\left(-\frac{1}{2}\right)^n
\]
当 \(n\in\N^*\) 时，\(-\frac{1}{2}\leqslant\left(-\frac{1}{2}\right)^n\leqslant\frac{1}{4}\)，所以 \(\frac{3}{4}\leqslant S_n\leqslant\frac{3}{2}\). 特别地 \(S_n>0\)，且
\[
\frac{13}{6}-S_n-\frac{1}{S_n}
=-\frac{(3S_n-2)(2S_n-3)}{6S_n}\geqslant 0
\]
因为 \(S_n\in\left[\frac{3}{4},\frac{3}{2}\right]\subset\left[\frac{2}{3},\frac{3}{2}\right]\). 因此
\(S_n+\frac{1}{S_n}\leqslant\frac{13}{6}\).
\end{solution}

\begin{problem}
设 \(a \in [-2,0]\)，已知函数 \(f(x) = \begin{cases}
x^3 - (a + 5)x, & x \leqslant 0,\\
x^3 - \frac{a + 3}{2}x^2 + ax, & x > 0.
\end{cases}\)
\begin{enumerate}
\item 证明： \( f(x) \) 在区间 \( (-1,1) \) 内单调递减，在区间 \( (1,+\infty) \) 内单调递增；
\item 设曲线 \( y = f(x) \) 在点 \( P_i(x_i, f(x_i)) \) （\(i = 1\)，\(2\)，\(3\)） 处的切线相互平行，且 \( x_{1}x_{2}x_{3} \neq 0 \)，证明：\( x_{1} + x_{2} + x_{3} > -\frac{1}{3} \).
\end{enumerate}
\end{problem}

\begin{solution}
记
\(f_1(x)=x^3-(a+5)x\)，\((x\leqslant0)\)，\(f_2(x)=x^3-\frac{a+3}{2}x^2+ax\)，\((x>0)\).
\begin{enumerate}
\item 当 \(-1<x<0\) 时，
\(f_1'(x)=3x^2-(a+5)<3-(a+5)\leqslant0\)，且不等号严格成立；所以 \(f_1\) 在 \((-1,0)\) 内单调递减.

当 \(x>0\) 时，
\[
f_2'(x)=3x^2-(a+3)x+a=(3x-a)(x-1)
\]
由于 \(a\leqslant0\)，当 \(0<x<1\) 时 \(3x-a>0\)，\(x-1<0\)，所以 \(f_2'(x)<0\)；当 \(x>1\) 时两个因子均为正，所以 \(f_2'(x)>0\). 又当 \(-1<x<0\) 时，\(f_1(x)>f_1(0)=0\)，当 \(0<x<1\) 时，\(f_2(x)<f_2(0)=0\)，故跨过 \(x=0\) 仍保持严格递减. 因此 \(f(x)\) 在 \((-1,1)\) 内单调递减，在 \((1,+\infty)\) 内单调递增.

\item 设三条切线的公共斜率为 \(m\). 因为 \(x_i\neq0\)，且
\(f_1''(x)=6x<0\)（\(x<0\)），所以方程 \(f_1'(x)=m\) 在
\((-\infty,0)\) 内至多有一个解；方程 \(f_2'(x)=m\) 是关于 \(x\) 的二次方程，在
\((0,+\infty)\) 内至多有两个解. 所以不妨设
\(x_1<0<x_2<x_3\). 其中 \(x_2\)，\(x_3\) 是方程
\[
3x^2-(a+3)x+a-m=0
\]
的两个正根，故
\(x_2+x_3=\frac{a+3}{3}\)，\(\frac{a-m}{3}=x_2x_3>0\).
于是 \(m<a\). 又 \(x_1\) 满足 \(3x_1^2-(a+5)=m\)，所以
\[
3x_1^2=m+a+5<2a+5
\]
从而
\[
x_1> -\sqrt{\frac{2a+5}{3}}
\]
对 \(a\in[-2,0]\)，有
\[
\frac{(a+4)^2}{9}-\frac{2a+5}{3}=\frac{(a+1)^2}{9}\geqslant0
\]
且 \(a+4>0\)，所以 \(\sqrt{\frac{2a+5}{3}}\leqslant\frac{a+4}{3}\). 因此
\[
x_1+x_2+x_3> -\frac{a+4}{3}+\frac{a+3}{3}=-\frac{1}{3}
\]
\end{enumerate}
\end{solution}
